% !TeX root = ../main.tex

% ==================================================
% A.6. 데코레이터의 정의와 @decorator 변환식
% ==================================================

\subsection{데코레이터의 정의와 \zsh{@decorator} 변환식}

Decorator · Function Object · Name Rebinding · Closure · Definition Time · Call Time

\subsubsection{이전 단원과의 연결}

\zsh{A.1}부터 \zsh{A.5}까지 학습한 내용은 데코레이터(Decorator)를 이해하기 위한 기반이었다.

\begin{dhktable}
|단원|	학습한 원리|	데코레이터와의 연결|
|:---|:---|:---|
|A.1|	함수도 객체이며 이름에 바인딩된다.|	함수 객체를 인수로 전달하고 다른 함수 객체를 반환할 수 있다.|
|A.2|	중첩 함수는 정의 위치에 따른 렉시컬 스코프를 가진다.|	데코레이터 내부에서 새로운 함수를 정의할 수 있다.|
|A.3|	이름의 스코프는 정적으로 판정된다.|	내부 함수는 외부 매개변수를 자유 변수로 참조할 수 있다.|
|A.4|	클로저는 외부 바인딩에 접근하는 환경을 유지한다.|	반환된 함수가 원본 함수 객체를 계속 참조할 수 있다.|
|A.5|	nonlocal은 외부 바인딩의 재바인딩을 허용한다.|	데코레이터가 호출 횟수 등의 상태를 유지할 수 있다.|
\end{dhktable}

이제 이 개념들을 하나의 실행 구조로 통합할 것이다.

다음 코드를 살펴보자.

\begin{dhkcode}
def decorator(func):

    def wrapper():
        print("Before")
        result = func()
        print("After")
        return result

    return wrapper


@decorator
def greet():
    print("Hello")
    return 10


print(greet())
\end{dhkcode}

예상 출력:

\begin{dhkzsh}
Before
Hello
After
10
\end{dhkzsh}

표면적으로는 \zsh{@decorator}라는 문법이 \zsh{greet()}의 실행 전후에 기능을 추가한 것처럼 보인다.

하지만 Python의 객체 모델에서는 훨씬 구체적인 과정이 발생한다.

\begin{dhksmry}
    1. 원본 greet 함수 객체가 생성된다.

    2. 그 함수 객체가 decorator에 전달된다.

    3. decorator가 새로운 wrapper 함수를 생성한다.

    4. wrapper는 클로저를 통해 원본 함수 객체를 참조한다.

    5. decorator가 반환한 wrapper 함수 객체가 
       greet 이라는 이름에 바인딩된다.

    6. 이후 greet()을 호출하면 원본 함수가 아니라 
       wrapper가 먼저 실행된다.
\end{dhksmry}

이번 학습에서 가장 중요한 질문은 다음과 같다.

\begin{dhksmry}[centered]
@decorator는 기존 함수 객체를 수정하는 문법일까, 아니면

함수 객체를 전달하고 반환값으로 이름을 다시 바인딩하는 문법일까?
\end{dhksmry}

정답은 일반적으로 두 번째다.

다만 데코레이터가 반드시 새로운 함수를 반환해야 하는 것은 아니라는 점까지 정확하게 이해해야 한다.

\subsubsection{데코레이터의 정확한 정의}

Python에서 데코레이터는 함수나 클래스를 정의할 때 적용하는 호출 가능한 객체(Callable)다.

함수 데코레이터의 기본 구조는 다음과 같다.

\begin{dhkcode}
def decorator(func):
    # func를 이용한 처리
    return replacement
\end{dhkcode}

여기서 \zsh{func}는 데코레이션 대상인 함수 객체를 참조하고,

\zsh{replacement}는 그 이름에 최종적으로 바인딩될 객체다.

일반적인 함수 기반 데코레이터에서는 \zsh{replacement}가 새로운 함수 객체인 경우가 많다.

하지만 Python 문법이 반드시 새로운 함수 객체를 반환하도록 강제하지는 않는다.

핵심은 다음 변환식이다.

\begin{dhknote}
\begin{dhksmry}[centered]
데코레이터의 의미론적 변환 - 데코레이터 문법
\end{dhksmry}
\end{dhknote}

\begin{dhkcode}
@decorator
def greet():
    return 10
\end{dhkcode}

\begin{dhknote}
\begin{dhksmry}[centered]
데코레이터의 의미론적 변환 - 대응하는 일반 함수 호출
\end{dhksmry}
\end{dhknote}

\begin{dhkcode}
def greet():
    return 10

greet = decorator(greet)
\end{dhkcode}

두 코드는 일반적인 단일 데코레이터 상황에서 같은 의미를 가진다.

여기서 다음 표현식을 정확하게 분석해야 한다.

\begin{dhkcode}
greet = decorator(greet)
\end{dhkcode}

오른쪽의 \zsh{greet}은 함수 객체를 참조하는 이름이다.

\zsh{decorator(greet)}은 그 함수 객체를 인수로 전달하는 호출 표현식이다.

그리고 호출 결과가 왼쪽 이름 \zsh{greet}에 다시 바인딩된다.

따라서 데코레이터의 본질은 다음과 같이 정리할 수 있다.

\begin{dhksmry}[centered]
데코레이터는 대상 함수 객체를 받아 어떤 객체를 반환하며,

데코레이터 문법은 그 반환 객체를 

원래 정의 대상의 이름에 바인딩하는 구조다.
\end{dhksmry}

여기서 데코레이터의 기능과 \zsh{@} 문법의 역할을 구분해야 한다.

\begin{dhksmry}
    - 데코레이터가 무엇을 반환할지는 데코레이터 구현에 따라 결정된다.

    - @ 문법은 데코레이터를 적용하고 그 결과를 정의 대상의 이름에
      바인딩하는 역할을 한다.
\end{dhksmry}

\subsubsection{데코레이터를 사용하지 않고 먼저 구현해 보자}

\zsh{@} 문법을 잠시 제외하고, 일반적인 함수 객체 전달과 반환만으로 동일한 구조를 만들어 보자.

\begin{dhkcode}
def decorator(func):
    print("decorator 실행")

    def wrapper():
        print("Before")

        result = func()

        print("After")

        return result

    return wrapper


def greet():
    print("Hello")
    return 10


original = greet

greet = decorator(greet)

print("정의 및 데코레이션 완료")

result = greet()

print("반환값:", result)
print("같은 함수 객체인가?", original is greet)
\end{dhkcode}

예상 결과:

\begin{dhkzsh}
decorator 실행
정의 및 데코레이션 완료
Before
Hello
After
반환값: 10
같은 함수 객체인가? False
\end{dhkzsh}

여기서 \zsh{original}이라는 이름을 추가한 이유는 원본 함수 객체와 데코레이션 후 함수 객체를 비교하기 위해서다.

이 코드를 단계별로 분석해 보자.

\begin{dhknote}
\begin{dhksmry}[centered]
1. decorator 함수 객체 생성
\end{dhksmry}
\end{dhknote}

\begin{dhkcode}
def decorator(func):
    ...
\end{dhkcode}

모듈 코드에서 \zsh{def decorator} 문장이 실행되면 함수 객체가 생성되고

\zsh{decorator}라는 이름에 바인딩된다.

이 시점에는 \zsh{decorator}의 본문이 실행되지 않는다.

\begin{dhknote}
\begin{dhksmry}[centered]
2. 원본 greet 함수 객체 생성
\end{dhksmry}
\end{dhknote}

\begin{dhkcode}
def greet():
    print("Hello")
    return 10
\end{dhkcode}

새로운 함수 객체가 생성되고 \zsh{greet}에 바인딩된다.

아직 \zsh{greet()}을 호출하지 않았으므로 \zsh{"Hello"}는 출력되지 않는다.

\begin{dhknote}
\begin{dhksmry}[centered]
3. 원본 함수 객체의 참조 보존
\end{dhksmry}
\end{dhknote}

\begin{dhkcode}
original = greet
\end{dhkcode}

\zsh{original}이라는 이름이 원본 함수 객체를 참조하게 된다.

새로운 함수 객체가 생성되는 것은 아니다.

이 시점에는 다음 관계가 성립한다.

\begin{dhkdiagram}[in=6]
node n1 (1,1) :: original
node n2 (3,1) :: greet
node n3 (2,3) :: 원본 greet 함수 객체

merge n1, n2 -> n3
\end{dhkdiagram}

\begin{dhknote}
\begin{dhksmry}[centered]
4. decorator(greet) 호출
\end{dhksmry}
\end{dhknote}

\begin{dhkcode}
greet = decorator(greet)
\end{dhkcode}

먼저 오른쪽 표현식을 평가한다.

\zsh{decorator}가 참조하는 함수 객체를 호출하고, 원본 \zsh{greet} 함수 객체를 인수로 전달한다.

따라서 \zsh{decorator} 호출의 지역 매개변수 \zsh{func}는 원본 함수 객체를 참조한다.

\begin{dhkdiagram}[in=6]
node n1 (1,1) :: 모듈 이름 greet
node n2 (2,1) :: 모듈 이름 original
node n3 (3,1) :: 매개변수 func
node n4 (2,3) :: 원본 greet 함수 객체

merge n1,n2,n3 ->  n4
\end{dhkdiagram}

이제 \zsh{decorator} 본문이 실행되므로 다음 문장이 출력된다.

\begin{dhkzsh}
decorator 실행
\end{dhkzsh}

\begin{dhknote}
\begin{dhksmry}[centered]
5. wrapper 함수 객체 생성
\end{dhksmry}
\end{dhknote}

다음 중첩 함수 정의문이 실행된다.

\begin{dhkcode}
def wrapper():
    print("Before")

    result = func()

    print("After")

    return result
\end{dhkcode}

새로운 \zsh{wrapper} 함수 객체가 생성된다.

여기서 \zsh{wrapper}는 외부 함수 \zsh{decorator}의 매개변수 \zsh{func}를 참조한다.

따라서 \zsh{func}는 \zsh{wrapper}의 자유 변수(Free Variable)다.

이것이 \zsh{A.4}에서 학습한 클로저의 형성 조건과 연결된다.

\begin{dhknote}
\begin{dhksmry}[centered]
6. wrapper 함수 객체 반환
\end{dhksmry}
\end{dhknote}

\begin{dhkcode}
return wrapper
\end{dhkcode}

\zsh{decorator}는 \zsh{wrapper} 함수 객체를 반환한다.

여기서 \zsh{wrapper()}를 호출한 것이 아니라, 함수 객체 자체를 반환했다는 점이 중요하다.

따라서 아직 \zsh{"Before"}와 \zsh{"After"}는 출력되지 않는다.

\begin{dhknote}
\begin{dhksmry}[centered]
7. greet 이름 재바인딩
\end{dhksmry}
\end{dhknote}

\zsh{decorator(greet)} 호출 결과로 반환된 \zsh{wrapper} 함수 객체가 왼쪽의 \zsh{greet}에 바인딩된다.

이제 객체 관계는 다음과 같다.

\begin{dhkgraph}[uniform]
node n1 (1,1) :: original
node n2 (1,2) [accent] :: 원본 함수 객체 // 실제 Hello 실행 코드
node n3 (2,2) :: Cell func // 클로저 셀
node n4 (3,1) :: greet
node n5 (3,2) [accent] :: wrapper 함수 객체 // Before / func() / After

explain :: 데코레이션 이후 객체 관계
caption :: greet은 wrapper를 참조하지만, 원본 함수 객체는 original과 wrapper의 클로저를 통해 계속 참조된다.

edge n1 -> n2
edge n3 -> n2
edge n4 -> n5
edge n5 -> n3 :: 참조
\end{dhkgraph}

따라서 다음 비교는 \zsh{False}다.

\begin{dhkcode}
original is greet
\end{dhkcode}

두 이름이 서로 다른 함수 객체를 참조하기 때문이다.

\begin{dhknote}
\begin{dhksmry}[centered]
8. greet() 호출
\end{dhksmry}
\end{dhknote}

마지막으로 다음 문장이 실행된다.

\begin{dhkcode}
result = greet()
\end{dhkcode}

현재 \zsh{greet}은 \zsh{wrapper} 함수 객체를 참조하고 있다.

따라서 \zsh{wrapper}가 먼저 실행된다.

\zsh{wrapper} 내부에서 \zsh{func()}를 호출할 때 비로소 원본 함수 객체가 실행되는 것이다.

이것이 데코레이터의 기본 호출 구조다.

\subsubsection{\zsh{@decorator}는 언제 실행될까?}

이제 일반 함수 호출을 \zsh{@} 문법으로 변경해 보자.

\begin{dhkcode}
def decorator(func):
    print("decorating:", func.__name__)

    def wrapper():
        print("Before")
        result = func()
        print("After")
        return result

    return wrapper


@decorator
def greet():
    print("Hello")
    return 10


print("정의문 실행 완료")

print(greet())
\end{dhkcode}

예상 출력:

\begin{dhkzsh}
decorating: greet
정의문 실행 완료
Before
Hello
After
10
\end{dhkzsh}

이 출력 순서는 데코레이터의 실행 시점을 설명한다.

여기에는 세 가지 시점이 존재한다.

\begin{dhksmry}
    1. 함수 정의문 실행 시점
       원본 함수 객체가 생성되고 데코레이터가 적용된다.

    2. 이름 재바인딩 시점
       데코레이터 반환값이 greet에 바인딩된다.

    3. greet() 호출 시점
       현재 greet에 바인딩된 wrapper가 실행되고 원본 함수를 호출한다.
\end{dhksmry}

여기서 다음 두 문장을 구별해야 한다.

\zsh{@decorator}가 적용되는 시점은 함수 정의문이 실행되는 시점이다.

\zsh{wrapper} 함수 본문이 실행되는 시점은 일반적으로 데코레이션 이후 해당 함수 객체가 호출되는 시점이다.

두 시점은 다르다.

또한 여기서 말하는 함수 정의 시점은 반드시 모듈을 처음 불러오는 시점만 의미하지 않는다.

예를 들어 중첩 함수에 데코레이터가 붙어 있다면, 그 중첩 함수의 \zsh{def} 문장이 실행될 때마다 데코레이터가 적용될 수 있다.

\begin{dhkcode}
def decorator(func):
    print("decorating")

    def wrapper():
        return func()

    return wrapper


def outer():

    @decorator
    def inner():
        return 10

    return inner


f1 = outer()
f2 = outer()
\end{dhkcode}

예상 출력:

\begin{dhkzsh}
decorating
decorating
\end{dhkzsh}

\zsh{outer()}를 두 번 호출했기 때문에 내부의 \zsh{def inner} 문장도 두 번 실행된다.

따라서 서로 다른 \zsh{inner} 함수 객체가 생성되고, 데코레이터도 각각 적용되는 것이다.

여기서 정확한 기준은 소스 코드가 어디에 적혀 있는가가 아니라 해당 정의문이 언제 실행되는가라는 사실이다.

\subsubsection{함수 정의, 데코레이션, 함수 호출을 구별하기}

세 시점을 더 명확하게 확인해 보자.

\begin{dhkcode}
events = []


def decorator(func):
    events.append("2. decorator 호출")

    def wrapper():
        events.append("4. wrapper 호출")

        result = func()

        events.append("6. wrapper 반환")

        return result

    return wrapper


events.append("1. 정의문 실행 전")


@decorator
def greet():
    events.append("5. 원본 greet 호출")
    return "Hello"


events.append("3. 정의문 실행 후")

print(events)

result = greet()

print(result)
print(events)
\end{dhkcode}

예상 출력:

\begin{dhkzsh}
['1. 정의문 실행 전', '2. decorator 호출', '3. 정의문 실행 후']
Hello
['1. 정의문 실행 전', '2. decorator 호출', '3. 정의문 실행 후', 
'4. wrapper 호출', '5. 원본 greet 호출', '6. wrapper 반환']
\end{dhkzsh}

이 코드는 다음 사실을 확인해 준다.

\zsh{greet()}을 호출하기 전에도 \zsh{decorator}는 이미 실행되었다.

하지만 \zsh{wrapper} 본문과 원본 \zsh{greet} 본문은 아직 실행되지 않았다.

이것은 다음 세 호출을 혼동하면 안 된다는 의미다.

\begin{dhktable}
|표현식|	실제 의미|
|:---|:---|
|decorator(greet)|	원본 함수 객체를 인수로 전달하여 데코레이터 호출|
|greet()|	현재 greet에 바인딩된 객체 호출|
|func()|	wrapper 내부에서 자유 변수 func가 참조하는 원본 함수 호출|
\end{dhktable}

데코레이터가 함수 객체를 인수로 받는 것과 그 함수 객체를 호출하는 것은 전혀 다른 연산이다.

한 가지 경계 조건도 기억해야 한다.

데코레이터가 원본 함수를 자동으로 호출하는 것은 아니지만, 데코레이터 구현 안에 명시적으로 \zsh{func()}가 있다면 데코레이션 도중 원본 함수가 실행될 수도 있다.

따라서 원본 함수가 언제 실행되는지는 실제 데코레이터 구현을 확인해야 한다.

\subsubsection{\zsh{func}와 \zsh{func()} - 가장 중요한 구분}

다음 두 코드는 같은 의미가 아니다.

\begin{dhkcode}
greet = decorator(greet)

greet = decorator(greet())
\end{dhkcode}

첫 번째 코드에서는 오른쪽의 \zsh{greet}이 함수 객체를 참조한다.

따라서 데코레이터의 매개변수 \zsh{func}에는 함수 객체가 전달된다.

반면 두 번째 코드에서는 먼저 \zsh{greet()}이 실행된다.

그다음 그 반환값이 \zsh{decorator}에 전달된다.

실제 차이를 확인해 보자.

\begin{dhkcode}
def decorator(func):
    print("전달받은 타입:", type(func).__name__)

    def wrapper():
        return func()

    return wrapper


def greet():
    return 10


greet = decorator(greet())
\end{dhkcode}

예상 출력:

\begin{dhkzsh}
전달받은 타입: int
\end{dhkzsh}

여기서 중요한 점이 있다.

이 코드가 반드시 \zsh{decorator(greet())} 호출 시점에 오류를 발생시키는 것은 아니다.

위 예제의 \zsh{decorator}는 전달받은 객체가 호출 가능한지 즉시 검사하지 않는다.

단순히 그 객체를 클로저에 보관하는 \zsh{wrapper}를 생성해서 반환한다.

따라서 이후 다음 문장을 실행할 때

\begin{dhkcode}
greet()
\end{dhkcode}

\zsh{wrapper} 내부의 \zsh{func()}에서 오류가 발생한다.

\begin{dhkzsh}
TypeError: 'int' object is not callable
\end{dhkzsh}

왜냐하면 \zsh{func}가 함수 객체가 아니라 정수 객체 10을 참조하고 있기 때문이다.

이를 표로 정리하면 다음과 같다.

\begin{dhktable}
|표현식|	데코레이터에 전달되는 객체|
|:---|:---|
|decorator(greet)|	greet가 참조하는 함수 객체|
|decorator(greet())|	greet() 호출의 반환 객체|
|decorator(lambda: 10)|	새로 생성된 함수 객체|
|decorator(10)|	정수 객체 10|
\end{dhktable}

데코레이터의 매개변수 이름이 \zsh{func}라고 해서 Python이 자동으로 함수 타입을 보장하는 것은 아니다.

전달받은 객체를 어떻게 사용할지는 데코레이터 구현에 달려 있다.

\subsubsection{원본 함수 객체는 어디에 보관될까?}

이제 \zsh{A.4}와 \zsh{A.5}에서 학습한 클로저 구조를 직접 연결해 보자.

다음 코드를 실행해 보자.

\begin{dhkcode}
def decorator(func):

    def wrapper():
        return func()

    return wrapper


@decorator
def greet():
    return 10


print(greet.__name__)
print(greet.__code__.co_freevars)
print(greet.__closure__[0].cell_contents.__name__)
\end{dhkcode}

예상 출력:

\begin{dhkzsh}
wrapper
('func',)
greet
\end{dhkzsh}

이 출력은 세 가지 사실을 보여준다.

\begin{dhknote}
\begin{dhksmry}[centered]
1. 현재 greet은 wrapper 함수 객체를 참조한다
\end{dhksmry}
\end{dhknote}

\begin{dhkcode}
print(greet.__name__)
\end{dhkcode}

결과는 \zsh{"wrapper"}다.

왜냐하면 \zsh{greet}이라는 이름이 데코레이터의 반환값인 \zsh{wrapper} 함수 객체에 재바인딩되었기 때문이다.

\zsh{greet}이라는 이름이 붙어 있다고 해서 그 이름이 원본 함수 객체를 참조한다는 뜻은 아니다.

\begin{dhknote}
\begin{dhksmry}[centered]
2. wrapper는 func를 자유 변수로 사용한다
\end{dhksmry}
\end{dhknote}

\begin{dhkcode}
print(greet.__code__.co_freevars)
\end{dhkcode}

결과는 다음과 같다.

\begin{dhkzsh}
('func',)
\end{dhkzsh}

\zsh{wrapper} 본문은 외부 함수 \zsh{decorator}의 매개변수 \zsh{func}를 참조한다.

따라서 \zsh{func}는 자유 변수다.

\begin{dhknote}
\begin{dhksmry}[centered]
3. 클로저 셀은 원본 함수 객체를 참조한다
\end{dhksmry}
\end{dhknote}

\begin{dhkcode}
greet.__closure__[0].cell_contents
\end{dhkcode}

이 표현식은 \zsh{wrapper}의 클로저 셀에 바인딩된 객체를 반환한다.

이 예제에서는 그 객체가 원본 \zsh{greet} 함수 객체다.

따라서 다음 코드는 \zsh{False}를 출력한다.

\begin{dhkcode}
original = greet.__closure__[0].cell_contents

print(original is greet)
print(original.__name__)
\end{dhkcode}

예상 출력:

\begin{dhkzsh}
False
greet
\end{dhkzsh}

현재 \zsh{greet}이 참조하는 함수 객체와 \zsh{original}이 참조하는 함수 객체가 서로 다르기 때문이다.

\begin{dhknote}
\begin{dhksmry}[centered]
객체 관계 정리
\end{dhksmry}
\end{dhknote}

\begin{dhkbranch}[in=8]
node main :: 모듈 네임 스페이스
handoff main -> a :: greet
arrow a
node a :: wrapper 함수 객체
branchline a -> a1 :: __code__
handoff a1 -> a2 :: co_freevars = ('func',)
stop a2
handoff a -> b :: __closure__
arrow b
node b :: Cell func
arrow b
node b :: 원본 greet 함수 객체
handoff b -> c :: __code__
arrow c
node c :: 원본 함수 코드
\end{dhkbranch}

여기서 \zsh{A.4}에서 학습한 원리가 그대로 적용된다.

\zsh{decorator()} 호출은 종료되었지만, 반환된 \zsh{wrapper} 함수 객체가 클로저 셀을 참조하고 있기 때문에 원본 함수 객체에 계속 접근할 수 있는 것이다.

단, 모든 데코레이터가 원본 함수를 클로저에 보관하는 것은 아니다.

이것은 지금 작성한 wrapper 기반 데코레이터의 구조다.

\subsubsection{데코레이터는 반드시 \zsh{wrapper}를 만들어야 할까?}

그렇지 않다.

데코레이터는 대상 함수 객체를 전달받아 다른 객체를 반환하는 구조이지, 반드시 중첩 함수를 생성해야 하는 것은 아니다.

다음 예제를 살펴보자.

\begin{dhkcode}
def identity(func):
    print("데코레이터 적용")
    return func


@identity
def greet():
    return "Hello"


print(greet())
print(greet.__name__)
print(greet.__closure__)
\end{dhkcode}

예상 출력:

\begin{dhkzsh}
데코레이터 적용
Hello
greet
None
\end{dhkzsh}

이 데코레이터는 새로운 \zsh{wrapper}를 생성하지 않는다.

전달받은 원본 함수 객체를 그대로 반환한다.

따라서 데코레이션 후에도 \zsh{greet}은 원본 함수 객체를 참조한다.

이 예제에서는 외부 함수의 바인딩을 캡처하지 않으므로 \zsh{__closure__}도 \zsh{None}이다.

즉, 데코레이터 적용과 클로저 생성은 반드시 함께 발생하는 것이 아니다.

\begin{dhknote}
\begin{dhksmry}[centered]
반환값이 함수가 아닌 경우
\end{dhksmry}
\end{dhknote}

다음 코드도 문법적으로 가능하다.

\begin{dhkcode}
def replace(func):
    return 100


@replace
def greet():
    return "Hello"


print(greet)
print(type(greet))
\end{dhkcode}

예상 출력:

\begin{dhkzsh}
100
<class 'int'>
\end{dhkzsh}

여기서는 \zsh{greet}이 정수 객체 100에 바인딩된다.

따라서 이후 \zsh{greet()}을 호출하면 \zsh{TypeError}가 발생한다.

이 예제가 보여주는 원리는 중요하다.

데코레이터의 반환값이 반드시 함수 객체여야 한다는 언어 차원의 제약은 없다.

다만 데코레이션 이후에도 함수처럼 호출하려면 반환 객체가 호출 가능해야 한다.

실무에서 일반적인 함수 데코레이터가 \zsh{wrapper} 함수 객체를 반환하는 이유도 기존 함수의 호출 인터페이스를 유지하면서 동작을 확장하기 위해서다.

\subsubsection{이름 재바인딩은 언제 완료될까?}

앞에서는 데코레이터의 반환값이 원래 이름에 다시 바인딩된다고 설명했다.

그런데 그 바인딩은 데코레이터가 호출되기 전에 완료될까, 아니면 호출이 끝난 뒤 완료될까?

다음 코드를 살펴보자.

\begin{dhkcode}
def greet():
    return "Old"


def decorator(func):
    print("데코레이션 중:", greet())
    return func


@decorator
def greet():
    return "New"


print("데코레이션 후:", greet())
\end{dhkcode}

예상 출력:

\begin{dhkzsh}
데코레이션 중: Old
데코레이션 후: New
\end{dhkzsh}

왜 이런 결과가 발생할까?

처음에는 \zsh{greet}이라는 이름은 \zsh{"Old"}를 반환하는 함수 객체를 참조한다.

그 다음 새로운 \zsh{greet} 정의문이 실행된다.

이때 새로운 함수 객체가 생성되고 \zsh{decorator}에 전달된다.

하지만 아직 새로운 함수 객체가 모듈 이름 \zsh{greet}에 최종 바인딩되기 전이다.

따라서 \zsh{decorator} 내부에서 전역 이름 \zsh{greet}을 조회하면 기존 함수 객체를 찾게 된다.

이후 \zsh{decorator}가 새로운 함수 객체를 반환하고, 그 반환값이 \zsh{greet}에 바인딩된다.

그래서 마지막 호출에서는 \zsh{"New"}가 반환되는 것이다.

이 과정을 나타내면 다음과 같다.

\begin{dhkdiagram}[in=6]
node n1 (1,1) :: 초기 상태
node n2 (2,2) :: greet
node n3 (2,3) :: 기존 함수 객체

edge n1 -- n2 :: VH
edge n2 -> n3
\end{dhkdiagram}

\begin{dhkdiagram}[in=2]
node n3 (3,1) :: 새로운 def greet 실행
node n4 (4,2) :: 새로운 함수 객체 생성
node n5 (5,2) :: decorator(새 함수 객체)
node n6 (6,2) :: 반환 객체 획득
node n7 (7,2) :: greet 이름 재바인딩
node n8 (8,2) :: greet
node n9 (8,3) :: 반환 객체

edge n3 -- n4 :: VH
edge n4 -> n5
edge n5 -> n6
edge n6 -> n7
edge n7 -> n8
edge n8 -> n9
\end{dhkdiagram}

여기서 중요한 판단 기준은 다음이다.

데코레이터가 적용되는 동안 새로 정의한 함수 객체가 반드시 기존 이름에 먼저 바인딩되어 있는 것은 아니다.

최종 이름 바인딩은 데코레이터 적용 결과가 결정된 뒤 이루어진다.

\subsubsection{데코레이터 적용과 원본 함수의 내부 이름 탐색}

여기서 조금 더 깊이 생각해 보자.

데코레이션 후 원본 함수 객체는 클로저를 통해 유지될 수 있다.

그런데 원본 함수 본문에서 자기 자신의 이름을 참조한다면 어떤 일이 발생할까?

\begin{dhkcode}
def decorator(func):

    def wrapper(n):
        print("wrapper:", n)
        return func(n)

    return wrapper


@decorator
def countdown(n):
    if n == 0:
        return 0

    return countdown(n - 1)


print(countdown(2))
\end{dhkcode}

예상 출력:

\begin{dhkzsh}
wrapper: 2
wrapper: 1
wrapper: 0
0
\end{dhkzsh}

왜 \zsh{wrapper}가 세 번 실행될까?

원본 \zsh{countdown} 함수 본문에 다음 표현식이 있다.

\begin{dhkcode}
countdown(n - 1)
\end{dhkcode}

이 \zsh{countdown}은 함수 객체 자체를 직접 가리키는 특별한 자기 참조가 아니다.

일반적인 Python 이름 탐색 규칙에 따라 조회되는 전역 이름이다.

그런데 데코레이션 이후 모듈의 \zsh{countdown}은 원본 함수가 아니라 \zsh{wrapper}를 참조하고 있다.

따라서 원본 함수가 재귀 호출을 수행할 때도 \zsh{wrapper}를 다시 호출하게 되는 것이다.

실행 흐름을 살펴보자.

\begin{dhkbranch}[in=8]
node main :: countdown(2)
arrow main
node main :: wrapper(2)
arrow main
node main :: 원본 countdown(2)
handoff main -> a :: countdown(1)
arrow a
node a :: wrapper(1)
arrow a
node a :: 원본 countdown(1)
handoff a -> b :: countdown(0)
arrow b
node b :: wrapper(0)
arrow b
node b :: 원본 countdown(0)
arrow b
node b :: 0
\end{dhkbranch}

이 예제는 두 가지 원리를 연결한다.

첫째, 원본 함수 객체는 클로저에 의해 유지될 수 있다.

둘째, 원본 함수 본문에 작성된 전역 이름은 함수 정의 당시의 객체를 영구적으로 가리키는 것이 아니라, 실행 시점의 전역 바인딩을 조회한다.

따라서 데코레이터에 의한 이름 재바인딩은 원본 함수 내부의 재귀 호출 경로에도 영향을 줄 수 있다.

이것은 \zsh{A.3}의 이름 탐색 원리가 데코레이터와 결합된 사례다.

\subsubsection{\zsh{@decorator}와 \zsh{@decorator()}는 다르다}

다음 두 문법을 비교해 보자.

\begin{dhkcode}
@decorator
def greet():
    return 10
\end{dhkcode}

\begin{dhkcode}
@decorator()
def greet():
    return 10
\end{dhkcode}

두 표현식은 같은 의미가 아니다.

첫 번째는 \zsh{decorator}가 참조하는 호출 가능한 객체를 데코레이터로 사용한다.

따라서 개념적으로 다음과 대응한다.

\begin{dhkcode}
greet = decorator(greet)
\end{dhkcode}

반면 두 번째는 먼저 \zsh{decorator()}를 호출한다.

그리고 그 호출 결과로 얻은 객체를 실제 데코레이터로 사용한다.

개념적으로 다음과 대응한다.

\begin{dhkcode}
greet = decorator()(greet)
\end{dhkcode}

따라서 \zsh{decorator()}가 반환하는 객체도 호출 가능해야 한다.

이러한 구조는 인수를 받는 데코레이터에서 자주 사용된다.

예를 들어 다음과 같은 형태다.

\begin{dhkcode}
@repeat(3)
def greet():
    print("Hello")
\end{dhkcode}

이때 \zsh{repeat(3)}은 실제 데코레이터를 반환하는 데코레이터 생성기(Decorator Factory) 역할을 할 수 있다.

이 구조는 \zsh{A.9}에서 상세하게 학습할 것이다.

지금은 다음 구분만 정확하게 확립하면 된다.

\begin{dhktable}
|문법|	적용 방식|
|:---|:---|
|@decorator|	decorator 자체를 데코레이터로 사용|
|@decorator()|	decorator() 호출 결과를 데코레이터로 사용|
|@decorator(arg)|	decorator(arg) 호출 결과를 데코레이터로 사용|
\end{dhktable}

특히 \zsh{@decorator}를 \zsh{@decorator()}로 임의로 변경하면

함수 정의 시점에 \zsh{TypeError}가 발생하거나 전혀 다른 동작이 나타날 수 있다.

\subsubsection{wrapper의 인수 전달 문제}

지금까지 사용한 wrapper는 다음과 같았다.

\begin{dhkcode}
def wrapper():
    return func()
\end{dhkcode}

이 구조에는 한 가지 제한이 있다.

원본 함수가 인수를 받는다면 \zsh{wrapper}도 그 호출을 처리할 수 있어야 한다.

다음 코드를 살펴보자.

\begin{dhkcode}
def decorator(func):

    def wrapper():
        return func()

    return wrapper


@decorator
def add(a, b):
    return a + b


print(add(10, 20))
\end{dhkcode}

예상 결과:

\begin{dhkzsh}
TypeError: ... wrapper() takes 0 positional arguments 
but 2 were given
\end{dhkzsh}

왜 오류가 발생할까?

데코레이션 후 \zsh{add}는 원본 함수가 아니라 \zsh{wrapper}를 참조한다.

따라서 다음 호출은

\begin{dhkcode}
add(10, 20)
\end{dhkcode}

실제로는 \zsh{wrapper(10, 20)} 호출에 해당한다.

그런데 \zsh{wrapper}에는 매개변수가 없다.

따라서 원본 \zsh{add}에 도달하기 전에 \zsh{wrapper} 호출 단계에서 오류가 발생한다.

이 문제를 해결하는 대표적인 방법은 다음과 같다.

\begin{dhkcode}
def decorator(func):

    def wrapper(*args, **kwargs):
        return func(*args, **kwargs)

    return wrapper
\end{dhkcode}

여기서 \zsh{*args}와 \zsh{**kwargs}를 통해 \zsh{wrapper}가 받은 인수를 원본 함수에 전달할 수 있다.

하지만 단순히 인수를 전달하는 것만으로 모든 함수 호출 계약이 자동으로 보존되는 것은 아니다.

원본 함수의 반환값, 예외, 시그니처(Signature), 동기·비동기 동작 등을 어떻게 처리할지도 고려해야 한다.

이 부분은 다음 \zsh{A.7}에서 자세히 분석할 것이다.

\subsubsection{데코레이터 적용 후 함수 메타데이터는 어떻게 될까?}

앞에서 다음 결과를 확인했다.

\begin{dhkcode}
print(greet.__name__)
\end{dhkcode}

데코레이션 후 \zsh{greet}이 \zsh{wrapper} 함수 객체를 참조한다면 결과는 일반적으로 다음과 같다.

\begin{dhkzsh}
wrapper
\end{dhkzsh}

왜 그럴까?

함수 객체의 \zsh{__name__}은 그 함수 객체에 관한 메타데이터다.

모듈 이름 \zsh{greet}이 \zsh{wrapper}를 참조하도록 바뀌었다고 해서 \zsh{wrapper}함수 객체의 \zsh{__name__}이 자동으로 \zsh{"greet"}으로 변경되는 것은 아니다.

다음 코드를 살펴보자.

\begin{dhkcode}
def decorator(func):

    def wrapper():
        return func()

    return wrapper


@decorator
def greet():
    """인사말을 반환한다."""
    return "Hello"


print(greet.__name__)
print(greet.__doc__)
\end{dhkcode}

예상 출력:

\begin{dhkzsh}
wrapper
None
\end{dhkzsh}

원본 함수의 \zsh{__name__}은 \zsh{"greet"}이고, \zsh{__doc__}에는 독스트링이 존재한다.

하지만 데코레이션 후 \zsh{greet}은 새로운 \zsh{wrapper} 함수 객체를 참조하고 있다.

따라서 \zsh{wrapper}의 메타데이터가 나타나는 것이다.

이 문제를 해결하기 위해 일반적으로 \zsh{functools.wraps}를 사용한다.

\begin{dhkcode}
from functools import wraps


def decorator(func):

    @wraps(func)
    def wrapper():
        return func()

    return wrapper
\end{dhkcode}

다만 \zsh{functools.wraps}가 어떤 속성을 복사하고, \zsh{__wrapped__}를 어떻게 설정하며, \zsh{inspect.signature()}가 이를 어떻게 활용하는지는 \zsh{A.8}에서 다룰 것이다.

이번 학습에서는 이름 재바인딩이 원본 함수의 메타데이터를 자동으로 새 함수 객체에 복사하지 않는다는 점을 이해하면 충분하다.

\subsubsection{데코레이터의 내부 구조를 종합적으로 검증하기}

이제 지금까지의 내용을 하나의 독립적인 코드로 통합해 보자.

\begin{dhkcode}
def decorator(func):
    print("[1] decorator 호출")

    def wrapper():
        print("[3] wrapper 호출")

        result = func()

        print("[5] wrapper 종료")

        return result

    return wrapper


@decorator
def greet():
    print("[4] 원본 greet 호출")
    return 10


print("[2] 데코레이션 완료")

wrapped = greet

freevars = wrapped.__code__.co_freevars
cells = wrapped.__closure__

closure_map = dict(zip(freevars, cells))

original = closure_map["func"].cell_contents

print("현재 이름의 함수:", wrapped.__name__)
print("원본 함수:", original.__name__)
print("함수 객체 동일성:", wrapped is original)

print("반환값:", greet())
\end{dhkcode}

예상 실행 결과:

\begin{dhkzsh}
[1] decorator 호출
[2] 데코레이션 완료
현재 이름의 함수: wrapper
원본 함수: greet
함수 객체 동일성: False
[3] wrapper 호출
[4] 원본 greet 호출
[5] wrapper 종료
반환값: 10
\end{dhkzsh}

이 코드는 다음 다섯 가지 관계를 직접 확인한다.

\begin{dhknote}
\begin{dhksmry}[centered]
첫째, 데코레이터 호출은 함수 정의문 실행 중 발생한다
\end{dhksmry}
\end{dhknote}

따라서 [1]이 [2]보다 먼저 출력된다.

\begin{dhknote}
\begin{dhksmry}[centered]
둘째, 데코레이션 이후 greet은 wrapper를 참조한다
\end{dhksmry}
\end{dhknote}

따라서 \zsh{wrapped.__name__}은 \zsh{"wrapper"}다.

\begin{dhknote}
\begin{dhksmry}[centered]
셋째, 원본 함수 객체는 wrapper의 클로저를 통해 접근할 수 있다
\end{dhksmry}
\end{dhknote}

\zsh{closure_map["func"].cell_contents}는 원본 함수 객체를 반환한다.

\begin{dhknote}
\begin{dhksmry}[centered]
넷째, wrapper와 원본 함수는 서로 다른 객체다
\end{dhksmry}
\end{dhknote}

따라서 \zsh{wrapped is original}은 \zsh{False}다.

\begin{dhknote}
\begin{dhksmry}[centered]
다섯째, wrapper가 원본 함수를 호출하고 반환값을 전달한다
\end{dhksmry}
\end{dhknote}

따라서 [3], [4], [5] 순서로 출력되고 최종적으로 10이 반환된다.

여기서 \zsh{closure_map}을 사용한 이유도 중요하다.

\zsh{__closure__[0]}이 항상 원본 함수 객체를 참조한다고 일반화하면 안 된다.

\zsh{wrapper}가 여러 자유 변수를 사용한다면 셀이 여러 개 존재할 수 있다.

따라서 내부 구조를 분석할 때는 \zsh{co_freevars}의 이름과 \zsh{__closure__}의 셀을 대응시키는 것이 더 정확하다.

\subsubsection{기존 Python 개념과의 연결}

이번 학습에서 등장한 데코레이터는 Python의 별도 객체 모델을 새로 도입한 것이 아니다.

일반적인 객체 모델을 결합한 구조다.

\begin{dhktable}
|기존 개념|	데코레이터에서의 역할|
|:---|:---|
|First-class Function|	원본 함수 객체를 다른 함수의 인수로 전달|
|함수 객체 생성|	wrapper 함수 객체 생성|
|이름 바인딩|	데코레이터 반환 객체를 원래 이름에 바인딩|
|Lexical Scope|	wrapper가 외부 함수의 매개변수를 참조|
|Free Variable|	wrapper의 func|
|Closure Cell|	원본 함수 객체에 접근하는 바인딩 유지|
|Callable|	wrapper를 함수처럼 호출|
|반환|	wrapper가 원본 함수 호출 결과를 반환|
|Identity|	원본 함수와 wrapper가 서로 다른 객체인지 비교|
\end{dhktable}

특히 클래스의 메서드 바인딩과도 구별해야 한다.

\begin{dhkcode}
obj.method
\end{dhkcode}

에서 발생하는 메서드 바인딩은 디스크립터(Descriptor)와 속성 탐색 규칙에 연결된다.

반면

\begin{dhkcode}
@decorator
def method(self):
    ...
\end{dhkcode}

에서의 데코레이터 적용은 클래스 본문에서 함수 정의문이 실행될 때 발생하는 객체 변환 과정이다.

이후 그 객체가 인스턴스 속성으로 접근될 때는 다시 디스크립터와 메서드 바인딩 규칙이 관여할 수 있다.

두 메커니즘은 서로 다른 단계에서 작동한다.

이 관계는 \zsh{A.11}에서 자세히 분석할 것이다.

\subsubsection{실전 응용 - 함수 실행 기록}

데코레이터의 가장 대표적인 활용 중 하나는 원본 함수의 핵심 로직을 변경하지 않고 부가 기능을 추가하는 것이다.

예를 들어 다음과 같은 실행 기록 기능을 생각해 볼 수 있다.

\begin{dhkcode}
def log_call(func):

    def wrapper(*args, **kwargs):
        print(f"{func.__name__} 호출 시작")

        result = func(*args, **kwargs)

        print(f"{func.__name__} 호출 종료")

        return result

    return wrapper


@log_call
def multiply(a, b):
    return a * b


print(multiply(3, 4))
\end{dhkcode}

예상 출력:

\begin{dhkzsh}
multiply 호출 시작
multiply 호출 종료
12
\end{dhkzsh}

여기서 원본 \zsh{multiply} 함수는 곱셈 연산만 담당한다.

실행 기록 기능은 \zsh{wrapper}가 담당한다.

이 구조가 유용한 이유는 기능의 책임을 분리할 수 있기 때문이다.

원본 함수는 자신의 계산을 수행하고, 데코레이터는 호출 전후의 공통 처리를 수행한다.

실제 프로그램에서는 이러한 구조를 이용해 로깅(Logging), 권한 검사, 호출 횟수 기록, 실행 시간 측정, 캐싱(Caching) 등을 구현할 수 있다.

하지만 각각의 기능은 단순히 \zsh{@}를 붙이는 것만으로 자동 궇녀되는 것이 아니다.

데코레이터가 반환하는 객체가 실제로 어떤 동작을 수행하도록 작성되어 있는지가 핵심이다.

\subsubsection{자주 발생하는 오해와 경계 조건}

\begin{dhktable}
|오해|	정확한 설명|
|:---|:---|
|@decorator는 원본 함수 객체를 직접 수정한다.|	일반적으로 원본 객체를 전달하고 반환 객체를 이름에 다시 바인딩해.|
|@decorator는 decorator()와 같다.|	전자는 데코레이터 객체를 사용하고, 후자는 먼저 호출 결과를 얻어 사용해.|
|데코레이터에는 함수 호출 결과가 전달된다.|	일반적인 함수 데코레이터에는 함수 객체 자체가 전달돼.|
|wrapper는 데코레이션 시점에 실행된다.|	wrapper 함수 객체는 생성되지만 본문은 별도 호출이 있어야 실행돼.|
|데코레이터는 반드시 새로운 함수를 반환한다.|	원본 함수를 그대로 반환하거나 다른 종류의 객체를 반환할 수도 있어.|
|원본 함수 객체는 데코레이션 후 사라진다.|	wrapper의 클로저 등 다른 참조가 있다면 계속 유지될 수 있어.|
|데코레이션 후 함수 이름은 원본 함수 객체를 가리킨다.|	그 이름은 데코레이터가 반환한 객체를 가리켜.|
|wrapper가 원본 함수의 인수와 메타데이터를 자동으로 보존한다.|	인수 전달과 메타데이터 보존은 별도로 구현해야 해.|
\end{dhktable}

여기서 특히 강조할 점은 데코레이터 문법 자체와 특정 데코레이터의 구현을 구분해야 한다는 것이다.

\zsh{@decorator}가 적용되었다는 사실만으로 원본 함수의 호출 횟수, 반환값, 예외 처리 방식, 메타데이터가 어떻게 달라지는지 단정할 수 없다.

실제 동작은 데코레이터가 반환하는 객체와 그 객체의 구현에 따라 결정된다.

\subsubsection{핵심 정리와 판단 기준}

이번 학습에서 가장 중요한 변환식은 다음과 같다.

\[
\texttt{@decorator}
\quad\Longleftrightarrow\quad
\texttt{func = decorator(func)}
\]

이 표현은 단일 함수 데코레이터의 일반적인 의미를 나타낸 것이다.

실제 소스 코드가 내부적으로 반드시 이 대입문으로 문자 그대로 치환한다는 의미는 아니다.

이를 객체 관계로 해석하면 다음과 같다.

\begin{dhksmry}[centered]
원본 함수 정의
\da
원본 함수 객체 생성
\da
decorator(원본 함수 객체)
\da
데코레이터 반환 객체
\da
원래 함수 이름에 재바인딩
\da
이후 이름을 통한 호출
\end{dhksmry}

wrapper 기반 데코레이터라면 다음 구조가 추가된다.

\begin{dhkbranch}[in=8]
node main :: 함수 이름
arrow main
node main :: wrapper 함수 객체
branchline main -> a :: 호출 시 부가 기능 실행
stop a
handoff main -> c :: 클로저 셀
arrow c
node c :: 원본 함수 객체
arrow c
node c :: 원본 함수 호출
\end{dhkbranch}

이제 다음 질문에 스스로 답할 수 있다면 \zsh{A.6}의 핵심 구조를 이해한 것이다.

\begin{dhksmry}
    - @decorator는 언제 적용되는가?

    - 데코레이터에는 어떤 객체가 인수로 전달되는가?

    - 데코레이터가 반환한 객체는 어디에 바인딩되는가?

    - 원본 함수 객체와 wrapper 함수 객체는 동일한가?

    - wrapper가 원본 함수를 참조한다면 그 바인딩은 어디에 유지되는가?

    - 데코레이션 이후 func()를 호출하면 실제로 어떤 객체가 호출되는가?
\end{dhksmry}

가장 중요한 판단 기준은 하나다.

\begin{dhksmry}[centered]
데코레이터를 분석할 때는 함수 이름만 따라가지 말고,

각 시점에 그 이름이 어떤 객체를 참조하는지 추적해야 한다.
\end{dhksmry}